% SHARD-ID: AQM-41-CECH-COHOMOLOGY-QUANTUM-DESCENT
% SHARD-TITLE: Cech Cohomology and Quantum Descent
% SHARD-SUMMARY: Defines cover-relative Cech cohomology and the dual cosheaf homology used by the arithmetic-field net.
% SHARD-SUMMARY: Proves that the residue-label cosheaf has trivial higher Cech homology on finite covers.
% SHARD-SUMMARY: Identifies state, observable, and stabilizer descent defects as distinct quantum counterparts.
% SHARD-KEYWORDS: Cech cohomology, cosheaf homology, descent, marginal problem, stabilizer defect, observable net

\section{Cech Cohomology and Quantum Descent}
\label{sec:cech-cohomology-quantum-descent}

\subsection{Status and Source Anchors}

This shard adds the cover-complex layer requested after AQM-26--AQM-40.
Stacks Cohomology,
\path{references/algebraic_geometry/stacks_project_cohomology.tex}, lines
868--924, 926--946, 2438--2517, and 4071--4165, supplies the Cech complex,
\(H^0\) sheaf criterion, refinement colimit, and ordered Cech complex.  Curry,
\path{references/cosheaves/curry_1303_3255_source/ch_abstract_sheaves.tex},
lines 246--259 and 270--335, supplies the dual vector-space exactness and
Cech homology convention.  The quantum input is local: AQM-26 for labels and
observable nets, AQM-27 for states, and AQM-38--AQM-39 for stabilizer descent.
The new convention is \texttt{CONVENTIONS.md} (ae).

\subsection{Lamport Dependency Skeleton}

\[
\begin{array}{rcl}
D1&:&\mathcal U=\{U_i\}_{i\in I}\text{ is a finite ordered cover of }U.\\
D2&:&\check C^p(\mathcal U,F)=
      \prod_{i_0<\cdots<i_p}F(U_{i_0\cdots i_p}).\\
L1&:&d^2=0\text{ for the Cech differential}.\\
D3&:&\check H^p(\mathcal U,F)=\ker d^p/\operatorname{im}d^{p-1}.\\
L2&:&\check H^0(\mathcal U,F)\text{ is compatible local section data}.\\
D4&:&\check H^p(U,F)=\varinjlim_{\mathcal U}\check H^p(\mathcal U,F)
      \text{ when refinements are declared}.\\
D5&:&\check C_p(\mathcal U,G)=
      \bigoplus_{i_0<\cdots<i_p}G(U_{i_0\cdots i_p})
      \text{ for a precosheaf }G.\\
T1&:&\check H_0(\mathcal U,E)\cong E(U),\quad
      \check H_p(\mathcal U,E)=0\ (p>0).\\
D6&:&\check Z^0(\mathcal U,\mathsf S)
      \text{ is the set of compatible local states}.\\
T2&:&\mathsf S(U)\to\check Z^0(\mathcal U,\mathsf S)
      \text{ may fail injectivity and surjectivity}.\\
D7&:&L_{\leq U}(W)=L\cap E(W)\text{ for a global stabilizer }L\leq E(U).\\
T3&:&D_U(L;\mathcal U)\text{ is the cokernel of stabilizer Cech }H_0
      \text{ generation}.\\
P1&:&\mathcal A\text{ has an observable-net comparison map, not ordinary
      abelian Cech cohomology.}
\end{array}
\]

\subsection{Standard Cover Cohomology D1--D4}

Let \(X\) be a topological space and let
\(\mathcal U=\{U_i\}_{i\in I}\) be a finite open cover of \(U\subset X\).
Fix a total order on \(I\), and write
\[
  U_{i_0\cdots i_p}=U_{i_0}\cap\cdots\cap U_{i_p}.
\]
For an abelian presheaf \(F\), define
\[
  \check C^p(\mathcal U,F)=
  \prod_{i_0<\cdots<i_p}F(U_{i_0\cdots i_p}).
\]
If \(c\in\check C^p(\mathcal U,F)\), set
\[
  (dc)_{i_0\cdots i_{p+1}}
  =
  \sum_{j=0}^{p+1}(-1)^j\,
  c_{i_0\cdots \widehat{i_j}\cdots i_{p+1}}
  \big|_{U_{i_0\cdots i_{p+1}}}.
\]

\paragraph{Lemma \(L1\).}
The differential satisfies \(d^2=0\).

\emph{Proof.}
Fix \(i_0<\cdots<i_{p+2}\).  In \((d^2c)_{i_0\cdots i_{p+2}}\), the term
obtained by deleting \(i_a\) and \(i_b\), with \(a<b\), appears exactly twice.
The two restriction maps are equal by functoriality of the presheaf.  The
corresponding signs are
\[
  (-1)^a(-1)^{b-1}
  \quad\text{and}\quad
  (-1)^b(-1)^a .
\]
Their sum is \((-1)^{a+b-1}+(-1)^{a+b}=0\).  Summing over all pairs
\((a,b)\) gives \(d^2c=0\).

Thus
\[
  \check H^p(\mathcal U,F)
  =
  \ker\bigl(d:\check C^p\to\check C^{p+1}\bigr)/
  \operatorname{im}\bigl(d:\check C^{p-1}\to\check C^p\bigr)
\]
is well-defined.

\paragraph{Lemma \(L2\).}
\(\check H^0(\mathcal U,F)\) is the group of compatible local sections:
\[
  \check H^0(\mathcal U,F)
  =
  \left\{(s_i)\in\prod_iF(U_i):
  s_i|_{U_i\cap U_j}=s_j|_{U_i\cap U_j}\text{ for all }i,j\right\}.
\]
If \(F\) is a sheaf, the restriction map \(F(U)\to\check H^0(\mathcal U,F)\)
is bijective.

\emph{Proof.}
Since \(\check C^{-1}=0\), \(\check H^0=\ker(d:\check C^0\to\check C^1)\).
The displayed formula for \(d\) says exactly that all pairwise restrictions
agree.  The second sentence is precisely the sheaf gluing axiom, as in
Stacks Cohomology, lines 926--946.

When a cover-independent group is needed, we follow Stacks Cohomology,
lines 2438--2517, and define
\[
  \check H^p(U,F)=
  \varinjlim_{\mathcal U}\check H^p(\mathcal U,F),
\]
over open covers ordered by refinement.  In the arithmetic-field questions
below, the fixed cover is part of the physical question, so we keep the cover
in the notation unless this direct limit is explicitly invoked.

\subsection{Dual Cech Homology for Labels D5--T1}

Let \(G\) be a covariant precosheaf of abelian groups or vector spaces.  Define
\[
  \check C_p(\mathcal U,G)=
  \bigoplus_{i_0<\cdots<i_p}G(U_{i_0\cdots i_p}).
\]
For \(s\in G(U_{i_0\cdots i_p})\), put
\[
  \partial s=
  \sum_{j=0}^p(-1)^j
  r_{U_{i_0\cdots i_p},\,U_{i_0\cdots\widehat{i_j}\cdots i_p}}(s),
\]
where \(r\) is the covariant extension map.  The same sign cancellation as in
Lemma \(L1\) gives \(\partial^2=0\), hence \(\check H_p(\mathcal U,G)\).

For the residue-qudit label layer, \(E(W)=\bigoplus_{x\in W}E_x\), with
extension by zero along inclusions.  Assume the cover \(\mathcal U\) is finite
and \(E(U)\) is the finite-support direct sum used in AQM-24--AQM-26.

\paragraph{Theorem \(T1\).}
For every finite cover \(\mathcal U\) of \(U\),
\[
  \check H_0(\mathcal U,E)\cong E(U),
  \qquad
  \check H_p(\mathcal U,E)=0\quad(p>0).
\]

\emph{Proof.}
For a point \(x\in U\), let \(I_x=\{i\in I:x\in U_i\}\).  This set is
nonempty because \(\mathcal U\) covers \(U\).  Since direct sums distribute
over the finite Cech chain groups,
\[
  \check C_p(\mathcal U,E)
  \cong
  \bigoplus_{x\in U}
  \left(
    \bigoplus_{i_0<\cdots<i_p,\ i_j\in I_x}E_x
  \right).
\]
The differential preserves the summand indexed by \(x\), because all
extension maps are the identity on the \(E_x\) component.  Therefore the
\(x\)-summand is \(E_x\) tensored with the simplicial chain complex of the
full simplex on the vertex set \(I_x\).

Choose \(a\in I_x\).  The cone operator \(h(\sigma)=[a,\sigma]\), with the
orientation sign determined by inserting \(a\) into the ordered simplex,
satisfies \(\partial h+h\partial=\operatorname{id}-\epsilon\), where
\(\epsilon\) is the augmentation to one copy of \(E_x\) in degree \(0\).  Thus
the full simplex has homology \(E_x\) in degree \(0\) and \(0\) in higher
degrees.  Taking the direct sum over \(x\) gives the displayed result.

So the full label cosheaf carries no higher cover obstruction: every label is
assembled pointwise from the cover, and all higher cycles are simplex
boundaries.

\subsection{Quantum Counterparts D6--P1}

\paragraph{States.}
For the state presheaf \(\mathsf S(W)=\operatorname{State}(\mathcal A(W))\),
define only the zeroth compatible-family object
\[
  \check Z^0(\mathcal U,\mathsf S)=
  \left\{(\omega_i)\in\prod_i\mathsf S(U_i):
  \omega_i|_{U_i\cap U_j}=\omega_j|_{U_i\cap U_j}\right\}.
\]
There is a natural restriction map
\[
  \rho_{\mathcal U}:\mathsf S(U)\longrightarrow
  \check Z^0(\mathcal U,\mathsf S).
\]

\paragraph{Theorem \(T2\).}
\(\rho_{\mathcal U}\) is not generally injective and not generally
surjective.

\emph{Proof.}
Non-injectivity is AQM-27 Theorem \(T2\): on
\(\Spec(\mathbb Z/6\mathbb Z)\), an entangled state and a product state have
the same one-factor restrictions.  Non-surjectivity is AQM-27 Theorem \(T3\):
the \(AB\) and \(BC\) Bell marginals agree on \(B\), but no \(ABC\) state has
both as marginals.

Thus the state-side quantum counterpart is the marginal comparison map.
Failure of injectivity measures hidden global correlations.  Failure of
surjectivity measures incompatible marginal data.  These are not cohomology
groups yet, because \(\mathsf S\) is a convex set-valued presheaf, not an
abelian presheaf.

\paragraph{Stabilizers.}
Let \(L\leq E(U)\) be a global stabilizer label subspace.  Define the
supported sublabel precosheaf on opens \(W\subset U\) by
\[
  L_{\leq U}(W)=L\cap E(W).
\]
The augmentation \(\check C_0(\mathcal U,L_{\leq U})\to L\) sends a tuple of
local stabilizer labels to its sum in \(L\).  Its image is
\[
  \operatorname{Loc}_U(L;\mathcal U)
  =
  \sum_i(L\cap E(U_i)).
\]
The boundary from \(L\cap E(U_i\cap U_j)\) maps to zero under this
augmentation, because the two copies are the same global label with opposite
signs.  Hence the augmentation factors through
\(\check H_0(\mathcal U,L_{\leq U})\).

\paragraph{Theorem \(T3\).}
The stabilizer descent defect of AQM-38 is the cokernel of this zeroth Cech
homology generation map:
\[
  \operatorname{coker}\bigl(
  \check H_0(\mathcal U,L_{\leq U})\to L
  \bigr)
  \cong
  L/\operatorname{Loc}_U(L;\mathcal U)
  =
  D_U(L;\mathcal U).
\]

\emph{Proof.}
The image of \(\check C_0\to L\) is exactly
\(\operatorname{Loc}_U(L;\mathcal U)\).  Passing through the quotient
\(\check H_0\) does not change the image, because boundaries already map to
zero.  The cokernel is therefore the displayed quotient.  AQM-38 Theorem
\(T2\) proves that its vanishing is equivalent to local generation of the
stabilizer labels; AQM-39 gives the two-qutrit nonzero-defect example.

\paragraph{Observable nets.}
For a finite cover, let \(\operatorname{colim}_{\mathcal U}^{*{\rm alg}}\)
denote the ordinary colimit of the cover diagram of unital \(*\)-algebras.
The observable net supplies a canonical comparison map
\[
  \pi_{\mathcal U}:
  \operatorname{colim}_{\mathcal U}^{*{\rm alg}}\mathcal A
  \longrightarrow
  \mathcal A(U).
\]
AQM-26 proves that this ordinary algebra colimit is not the physical gluing:
for disjoint pieces it is a free-product gluing, while the Weyl net glues by
commuting tensor-local subalgebras inside the ambient algebra.  Therefore the
observable quantum counterpart is the comparison map \(\pi_{\mathcal U}\):
its image is the additive generated algebra, and its kernel records the extra
locality and commutation relations imposed by the physical representation.
Without abelianization or another declared linearization, this is not an
ordinary Cech cohomology group.

\subsection{Conclusion}

Classical Cech cohomology is the abelian cover-complex layer.  In AQF it
splits into layered descent tests: label Cech homology, stabilizer cokernels,
state marginal comparison maps, and observable-net colimit comparison maps.
Only the first two are cohomology or homology groups at this stage.
